\begin{abstract}
Linear-cost recurrent denoisers make repeated scans of long contexts affordable,
but useful refinement also requires accurate conditional predictions.
We use established learning/factorization-error accounting to distinguish these
requirements and study a 455M-parameter bidirectional RWKV denoiser on exactly
enumerable binary posterior tasks. A completed within-lineage intervention
shows that balancing evidence positions during fixed-length training reduces
far-context conditional error by 6.051 nats, with a conditional 95\% interval
$[5.392,6.735]$. This result motivates a stronger prospective test: can
source-calibrated errors select a reveal schedule on unseen dependency classes?
The frozen experiment trains six fresh task-training lineages and compares
policies with identical scan count and zero discarded dependence, sealing their
predicted rankings before held-out evaluation. \emph{Revision status: the new
experiment is running; no transfer outcome or independent-lineage success is
claimed in this working draft.} The earlier implementation-specific cost
certificate remains supporting evidence, rather than an optimized deployment
frontier.
\end{abstract}

\section{Introduction}
Long-context generation must extract distant evidence and coordinate the answer
tokens that depend on it. Linear attention and recurrent sequence mixers reduce
the arithmetic required to process the input. Masked diffusion provides a
separate mechanism: later calls condition unrevealed answer positions on tokens
revealed earlier. Combining the two can make informative additional scans
affordable, provided the learned conditionals remain accurate.

The usual factorization-error account explains why a suitable reveal schedule
helps an oracle. It does not predict whether a trained model can use that
schedule on a new context. An observed policy gain close to an oracle dependence
gap is consequently a consistency check once both learned errors are small.
Our scientific question is stronger: can measured conditional errors on source
tasks predict which equally costly schedule will work better on unseen
dependency structures?

We separate retrospective mechanism evidence from prospective prediction.
The completed position intervention holds context length, examples and update
count fixed; it supports improved distant-evidence use within one selected
training lineage. The new experiment compares two public information bases of
an invertible binary code. Both two-round policies have zero factorization loss,
so their difference isolates learned error. A source-calibrated predictor makes
its choices before target inference. Splitting whole code-invariant classes
prevents the held-out set from being only a relabeling of training structures.

The intended contribution is an empirical test of transferable conditional-error
prediction, grounded in established divergence accounting. Neither the KL
identity nor recurrent diffusion is claimed as new. Six fresh task-training
lineages address selection sensitivity for the new task, and all terminal
outcomes are retained. The older quality--cost certificate is kept within its
restricted comparison set; establishing a practical frontier would additionally
require competent matched causal and optimized attention baselines.
